% ============================================================================
\chapter{Tester son code~: Vitest, Jest, Playwright}
\label{ch-js-tests}
% ============================================================================

\begin{objectif}
\begin{itemize}
  \item D\'ecider quoi tester (et quoi ne surtout pas tester).
  \item \'Ecrire des tests lisibles avec le sch\'ema Arrange-Act-Assert et un nommage qui documente.
  \item Ma\^itriser les mocks, les faux timers, et le mock de r\'eseau (msw).
  \item Tester l'interface (composants) et le parcours complet (Playwright).
  \item Brancher les tests dans une CI GitHub Actions qui refuse un mauvais commit.
\end{itemize}
\end{objectif}

\section{Pourquoi et quoi}

Un test n'est pas une certification de correction~: c'est un \emph{filet} qui te permet de
modifier le code sans peur. Le retour sur investissement vient des zones \`a r\`egles m\'etier,
pas des lignes de branchement trivial.

\begin{tabz}{lXll}
  \headrow \thd{\'Etage} \& \thd{Quoi} \& \thd{Vitesse} \& \thd{Part saine} \\
  \midrule
  Unitaires \& fonctions pures, r\`egles m\'etier, parseurs, validateurs \& ms \& 70 \% \\
  Int\'egration \& un module + sa base/son disque/son API simul\'ee \& dizaines de ms \& 20 \% \\
  Bout-en-bout \& parcours utilisateur r\'eel dans un navigateur \& secondes \& 10 \% \\
\end{tabz}

Ne teste \textbf{pas}~: les getters/setters, une biblioth\`eque tiers (React ne t'int\'eresse pas),
le nombre exact d'appels internes d'une impl\'ementation (\c{c}a casse au moindre refactor), ni le CSS.

\section{Vitest~: le choix par d\'efaut en 2026}

\begin{terminal}
npm i -D vitest @vitest/coverage-v8 jsdom fake-indexeddb
npx vitest run --coverage
npx vitest --ui                 # interface de developpement : surveillance, filtres, timelines
\end{terminal}

\begin{codejson}
{
  "test": {
    "environment": "node",
    "globals": false,
    "coverage": {
      "provider": "v8",
      "reporter": ["text", "html", "lcov"],
      "thresholds": { "lines": 80, "functions": 75, "branches": 70 }
    },
    "projects": [
      { "extends": true, "test": { "include": ["tests/unit/**"], "environment": "node" } },
      { "extends": true, "test": { "include": ["tests/dom/**"], "environment": "jsdom" } }
    ]
  }
}
\end{codejson}

\begin{codejs}
// tests/unit/panier.test.js
import { describe, it, expect, beforeEach, vi } from "vitest";
import { calculerTotaux, ajouter } from "../../src/lib/panier.js";

describe("calculerTotaux", () => {
  it("applique la TVA de 20 % sur le montant hors taxes", () => {
    const t = calculerTotaux([{ prix: 10, qte: 2 }]);
    expect(t.ht).toBe(20);
    expect(t.tva).toBeCloseTo(4, 5);
    expect(t.ttc).toBeCloseTo(24, 5);
  });

  it.each([
    [[], 0],
    [[{ prix: 12.5, qte: 1 }], 15],
    [[{ prix: 8, qte: 3 }], 28.8],
  ])("panier %j -> TTC %f", (lignes, ttc) => {
    expect(calculerTotaux(lignes).ttc).toBeCloseTo(ttc * 1.2, 5);
  });

  it("refuse une quantite negative", () => {
    expect(() => ajouter({ lignes: [] }, { prix: 1, qte: -3 })).toThrow(/quantite/);
  });

  it("ne mute pas le panier recu", () => {
    const initial = Object.freeze({ lignes: Object.freeze([{ prix: 1, qte: 1 }]) });
    const suivant = ajouter(initial, { prix: 2, qte: 1 });
    expect(suivant.lignes).toHaveLength(2);
    expect(initial.lignes).toHaveLength(1);
  });
});
\end{codejs}

\subsection{Les assertions qui comptent}

\begin{codejs}
expect(valeur).toBe(1);              // identite (Object.is) : primitives
expect(objet).toEqual({ a: 1 });     // structure profonde, ignore undefined
expect(objet).toStrictEqual({ a: 1, b: undefined });   // strict : types + cles undefined
expect(tab).toEqual(expect.arrayContaining([2, 3]));
expect(objet).toMatchObject({ a: { b: 1 } });           // sous-ensemble : le test qui ne casse pas pour rien
expect(chaine).toMatch(/^\\d{4}-\\d{2}-\\d{2}$/);
expect(n).toBeCloseTo(0.3, 5);                           // flottants
expect(p).resolves.toEqual(x); expect(p).rejects.toThrow(Error);
expect(fn).toHaveBeenCalledTimes(2);
expect(fn).toHaveBeenCalledWith({ id: 1 }, expect.any(Number));
expect(element).toBeInTheDocument();                      // avec @testing-library/jest-dom
expect(localStorage.getItem("k")).not.toBeNull();
\end{codejs}

\section{Mocks, espions, horloges}

\begin{codejs}
import { describe, it, expect, vi, beforeEach, afterEach } from "vitest";

// 1. Espionner une fonction locale : on verifie le dialogue avec une dependance
const logger = { info: vi.fn() };
await service.traiter(logger);
expect(logger.info).toHaveBeenCalledWith(expect.stringContaining("ok"));
expect(logger.info.mock.calls.length).toBe(1);
logger.info.mockImplementation(() => { throw new Error("panne"); });
logger.info.mockResolvedValue({ id: 1 });       // pour une fonction asynchrone
logger.info.mockRestore();

// 2. Mock d'un module entier (ex : la couche reseau)
vi.mock("../../src/api.js", () => ({
  api: { get: vi.fn().mockResolvedValue({ items: [{ id: 1, titre: "x" }] }) },
}));

// 3. Horloges fausses : tester un debounce, un retry, une expiration de session
vi.useFakeTimers();
const fn = vi.fn();
debounce(fn, 300)("a"); debounce(fn, 300)("ab");
await vi.advanceTimersByTimeAsync(300);
expect(fn).toHaveBeenCalledTimes(1);
expect(fn).toHaveBeenCalledWith("ab");
vi.setSystemTime(new Date("2026-01-15T10:00:00Z"));
vi.useRealTimers();

// 4. Substituts d'environnement
process.env.NODE_ENV = "test";
globalThis.fetch = vi.fn().mockResolvedValue(new Response(JSON.stringify({ ok: true }), { status: 200 }));
\end{codejs}

\begin{astuce}
Quand le mock devient plus long que le test, c'est que le code est mal d\'ecoup\'e~: la d\'ependance
est trop dure \`a isoler. Injecte-la (param\`etre, constructeur, contexte) au lieu de mocker le
module. Les bons tests viennent d'un bon d\'ecouplage, pas d'un bon moqueur.
\end{astuce}

\subsection{Mock de r\'eseau avec msw~: le test qui ne ment pas}

\begin{codejs}
// tests/msw/handlers.js
import { http, HttpResponse } from "msw";
export const handlers = [
  http.get("/api/articles", ({ request }) => {
    const url = new URL(request.url);
    const page = Number(url.searchParams.get("page") ?? 1);
    return HttpResponse.json({ items: [], page, total: 0 });
  }),
  http.post("/api/articles", async ({ request }) => {
    const corps = await request.json();
    if (!corps.titre) return HttpResponse.json({ message: "titre requis" }, { status: 422 });
    return HttpResponse.json({ id: 1, ...corps }, { status: 201 });
  }),
  http.get("/api/articles/:id", () => HttpResponse.error()),        // simule une panne reseau
];
// tests/setup.js
import { beforeAll, afterEach, afterAll } from "vitest";
import { setupServer } from "msw/node";
export const server = setupServer(...handlers);
beforeAll(() => server.listen({ onUnhandledRequest: "error" }));   // toute requete non prevue = echec
afterEach(() => server.resetHandlers());
afterAll(() => server.close());
\end{codejs}

\section{Tester l'interface~: par le comportement visible}

\begin{codejs}
// tests/dom/liste.test.jsx
import { render, screen, fireEvent, waitFor } from "@testing-library/react";
import userEvent from "@testing-library/user-event";
import { ListeCourses } from "../../src/ui/ListeCourses.jsx";

it("ajoute un article et le rend comme termine au clic", async () => {
  render(<ListeCourses />);
  const champ = screen.getByRole("textbox", { name: /article/i });     // les roles d'abord :
  await userEvent.type(champ, "Riz{Enter}");                            //  accessible = testable
  expect(screen.getByRole("listitem")).toHaveTextContent("Riz");
  await userEvent.click(screen.getByRole("checkbox"));
  expect(screen.getByRole("listitem")).toHaveClass("fait");
  await waitFor(() => expect(localStorage.getItem("courses")).toContain("Riz"));
});
\end{codejs}

R\`egle d'or~: teste ce que l'\emph{utilisateur} voit et fait, jamais l'impl\'ementation interne.
\verb|getByRole|, \verb|getByLabelText|, \verb|getByText| plut\^ot que \verb|getByTestId|~; et si
tu n'arrives pas \`a s\'electionner un \'el\'ement par son r\^ole, c'est un vrai bug d'accessibilit\'e
que tu viens de trouver.

\section{Playwright~: le test bout-en-bout qui vaut le d\'etour}

\begin{terminal}
npm init playwright@latest         # installe navigateurs + github action + exemple
npx playwright test --ui
npx playwright codegen http://localhost:5173   # enregistre tes clics en code
npx playwright show-report
\end{terminal}

\begin{codejs}
// tests/e2e/parcours.spec.ts
import { test, expect } from "@playwright/test";

test("un visiteur peut ajouter un article et le retrouve apres rechargement", async ({ page, context }) => {
  await page.goto("/");
  await page.getByLabel("Article").fill("Riz basmati");
  await page.getByRole("button", { name: "Ajouter" }).click();

  await expect(page.getByRole("listitem")).toHaveText(/Riz basmati/);
  await expect(page.getByText(/1 article/)).toBeVisible();

  await page.reload();
  await expect(page.getByRole("listitem")).toHaveText(/Riz basmati/);
});

test("une requete API lente n'affiche pas de doublon", async ({ page }) => {
  let appelee = 0;
  await page.route("**/api/articles", async (route) => {
    appelee++;
    await new Promise((r) => setTimeout(r, 800));
    route.fulfill({ json: { items: [], total: 0 } });
  });
  await page.goto("/articles");
  await page.getByRole("button", { name: /recharger/i }).click();
  await page.getByRole("button", { name: /recharger/i }).click();
  await expect(page.getByTestId("compteur-appels")).toHaveText("1");
});
\end{codejs}

\begin{note}
Playwright sert aussi \`a autre chose~: captures d'\'ecran automatiques pour ton README,
PDFs, g\'en\'eration de vignettes, tests multi-navigateurs et multi-\emph{viewports},
mocks r\'eseau au niveau HTTP (le plus fid\`ele qui existe). Un seul outil, quatre usages.
\end{note}

\section{Tests property-based~: chercher les cas que tu n'as pas imagin\'es}

\begin{codejs}
import fc from "fast-check";
it("le tri est une permutation ordonnee, pour tout tableau", () => {
  fc.assert(
    fc.property(fc.array(fc.integer()), (tab) => {
      const tri = trier(tab);
      expect(tri).toHaveLength(tab.length);
      expect([...tri].sort((a, b) => a - b)).toEqual(tri);
      expect(tri.reduce((s, x) => s + x, 0)).toBe(tab.reduce((s, x) => s + x, 0));
    }),
    { numRuns: 200 }
  );
});
it("serialiser/parser est un aller-retour", () => {
  fc.assert(fc.property(panierArbitraire, (p) => {
    expect(deverser(verser(p))).toEqual(p);
  }));
});
\end{codejs}

C'est l'arme id\'eale pour les parseurs, les slugifieurs, les codecs, les calculateurs de prix.

\section{Tests et int\'egration continue}

\begin{codeyml}
# .github/workflows/ci.yml
name: CI
on:
  push: { branches: [main] }
  pull_request:
concurrency: { group: ${{ github.workflow }}-${{ github.ref }}, cancel-in-progress: true }
jobs:
  qualite:
    runs-on: ubuntu-latest
    steps:
      - uses: actions/checkout@v4
      - uses: pnpm/action-setup@v4
      - uses: actions/setup-node@v4
        with: { node-version: 22, cache: pnpm }
      - run: pnpm install --frozen-lockfile
      - run: pnpm lint && pnpm typecheck
      - run: pnpm test -- --coverage
      - uses: actions/upload-artifact@v4
        with: { name: couverture, path: coverage/lcov-report }
  e2e:
    runs-on: ubuntu-latest
    needs: qualite
    steps:
      - uses: actions/checkout@v4
      - uses: actions/setup-node@v4
        with: { node-version: 22 }
      - run: npm ci && npm run build
      - uses: microsoft/playwright-github-action@v1
      - run: npx playwright install --with-deps chromium
      - run: npm run start & npx wait-on http://localhost:4173 && npm run test:e2e
      - uses: actions/upload-artifact@v4
        if: failure()
        with: { name: traces, path: test-results/ }
\end{codeyml}

\begin{exercice}
\begin{enumerate}
  \item \'Ecris une suite de 15 tests sur le module \verb|panier|~: cas nominaux, cas limites
        (vide, quantit\'es nulles, prix n\'egatifs, d\'ecimales), erreurs, immuabilit\'e.
  \item Ajoute un test de non-r\'egression~: reproduis un bug que tu as vraiment eu cette semaine,
        \'ecris le test qui \'echoue, puis corrige le code pour le faire passer.
  \item Mock le r\'eseau avec msw et teste la gestion d'erreur de ton client HTTP~: panne,
        401 (d\'econnexion provoqu\'ee), 500, r\'eponse lente avec annulation.
  \item Teste un composant UI avec Testing Library~: seulement par r\^oles et libell\'es. Note
        chaque fois que tu as \'et\'e tent\'e par un \verb|data-testid|~: c'\'etait une dette d'accessibilit\'e.
  \item \'Ecris deux tests Playwright~: parcours d'ajout/suppression, et gestion du mode hors ligne.
  \item Configure la couverture \`a seuil 80 \% et la CI GitHub Actions. V\'erifie qu'un commit
        rouge \'echoue bien~: si tu n'as pas vu l'action passer au rouge, ta CI est d\'ecorative.
\end{enumerate}
\end{exercice}

\begin{ressources}
\begin{itemize}
  \item Vitest --- \url{https://vitest.dev/guide/}
  \item Jest (si tu dois le maintenir) --- \url{https://jestjs.io/fr/docs/getting-started}
  \item Testing Library --- \url{https://testing-library.com/docs/react-testing-library/intro/}
  \item Playwright --- \url{https://playwright.dev/docs/intro}
  \item msw --- \url{https://mswjs.io/docs/} --- fast-check --- \url{https://fast-check.dev/}
  \item GitHub Actions, doc compl\`ete --- \url{https://docs.github.com/fr/actions}
\end{itemize}
\end{ressources}

\index{tests}\index{Vitest}\index{Playwright}\index{mock}
